% Generated draft: requires mathematical and source review.
\documentclass[11pt]{article}
\usepackage[margin=25mm]{geometry}
\usepackage{fontspec}
\setmainfont{Latin Modern Roman}
\usepackage{amsmath,amssymb,mathtools,unicode-math}
\setmathfont{Latin Modern Math}
\usepackage{enumitem,needspace,xurl,hyperref}
\setlength{\parindent}{0pt}
\setlength{\parskip}{5pt}
\setlength{\emergencystretch}{3em}
\newcommand{\N}{\mathbb{N}}
\newcommand{\Z}{\mathbb{Z}}
\newcommand{\Q}{\mathbb{Q}}
\newcommand{\R}{\mathbb{R}}
\newcommand{\C}{\mathbb{C}}
\newcommand{\F}{\mathbb{F}}
\newcommand{\A}{\mathbb{A}}
\newcommand{\PP}{\mathbb{P}}
\newcommand{\E}{\mathbb{E}}
\newcommand{\eps}{\varepsilon}
\newcommand{\dd}{\,\mathrm{d}}
\begin{document}
\section*{UnsolvedMath ID: 701413}
\textbf{Source problem code:} catalogue entry
\subsection{Definitions and mathematical statement}
For a group \(G\), a subgroup \(B\leq G\), and \(x\in G\), write \(B^x=x^{-1}Bx\). For subsets \(U,V\subseteq G\), write \(UV=\{uv:u\in U,\ v\in V\}\). A subgroup \(A\leq G\) is \emph{\(G\)-permutable in \(G\)} if
\[
\forall B\leq G\ \exists x\in G:\quad AB^x=B^xA.
\]
A subgroup \(A\leq G\) is \emph{hereditarily \(G\)-permutable in \(G\)} if it is \(E\)-permutable in every subgroup \(E\) of \(G\) containing \(A\); explicitly,
\[
\forall E\ (A\leq E\leq G)\ \forall B\leq E\ \exists x\in E:\quad AB^x=B^xA.
\]
The conjugating element may depend on \(B\), and, in the hereditary condition, on \(E\) as well.

The source asks which finite non-abelian simple groups \(G\) possess:
\begin{enumerate}
\item a non-trivial \(G\)-permutable subgroup;
\item a non-trivial hereditarily \(G\)-permutable subgroup.
\end{enumerate}
Here a finite non-abelian simple group is a finite non-abelian group whose only normal subgroups are \(\{1\}\) and \(G\).

\textbf{Unresolved qualification concerning properness.} The questions say ``non-trivial'' without explicitly saying ``proper.'' With the usual meaning \(A\neq\{1\}\), the permitted range is \(\{1\}<A\leq G\), which allows \(A=G\); this subgroup satisfies both displayed definitions. The supplied comments instead discuss proper subgroups. Thus the apparent intended classification is of finite non-abelian simple groups satisfying, respectively,
\[
\exists A\ (\{1\}<A<G)\ \forall B\leq G\ \exists x\in G:\quad AB^x=B^xA
\]
and
\[
\exists A\ (\{1\}<A<G)\ \forall E\ (A\leq E\leq G)\ \forall B\leq E\ \exists x\in E:\quad AB^x=B^xA.
\]
These proper-subgroup formulations are explicitly conditional interpretations, not a verified correction of the source.

\textbf{Supplied comments of 2025.} The record reports that, among sporadic groups, only \(J_1\) has a proper \(G\)-permutable subgroup, citing Galt and Tyutyanov (2022). It also reports that sporadic, alternating, and exceptional groups of Lie type have no proper hereditarily \(G\)-permutable subgroups, citing the three references listed below. These are supplied literature claims, not independently verified results. Their use of ``proper'' also requires a non-triviality qualification: \(\{1\}\) satisfies both definitions. No precise range or low-degree qualifications for the named group families are supplied. The registry labels the problem open; this is a supplied status label rather than an independently verified assessment.
\subsection{Short English statement}
Classify the finite non-abelian simple groups having a subgroup that permutes with a suitable conjugate of every subgroup, and classify those having such a subgroup with the same property inside every intermediate subgroup. The apparent intended subgroup range is \(\{1\}<A<G\), but properness is not explicit in the questions. The supplied 2025 partial-result comments and open status require source and mathematical review.
\subsection{Sources}
\begin{itemize}
\item Local UnsolvedMath registry, supplied complete source record: permanent ID \(701413\), exact source problem code \texttt{catalogue entry}, title \emph{Kourovka Notebook Problem 17.112}. The record gives upstream revision \texttt{d371ee917ac9647fbb06cd6fc529deec6f3dda82} and update timestamp \texttt{2026-10-16T14:00:00Z}. These are registry provenance data, not a literature-verification date.
\item \emph{The Kourovka Notebook (Unsolved Problems in Group Theory)}, original problem 17.112, 17th issue (2010); the supplied background also identifies the 21st issue edition (2026), arXiv:1401.0300v48. Proposed by A. N. Skiba and V. N. Tyutyanov in 2010. The collection description names E. I. Khukhro and V. D. Mazurov as editors. This authoritative-source attribution is supplied by the record and has not been independently checked.
\item A. A. Galt and V. N. Tyutyanov, \emph{Siberian Math. J.}, \textbf{63}, no. 4 (2022), 691--698. Cited in the supplied comments of 2025 for the sporadic-group assertions. No article title or link is supplied.
\item A. F. Vasilyev and V. N. Tyutyanov, \emph{Izv. Gomel. Gos. Univ. F. Skoriny}, \textbf{2012}, no. 5(74) (2012), 148--150 (Russian). Cited in the supplied comments of 2025 concerning hereditary permutability. No article title or link is supplied.
\item A. A. Galt and V. N. Tyutyanov, \emph{Siberian Math. J.}, \textbf{64}, no. 5 (2023), 1110--1116. Cited in the supplied comments of 2025 concerning hereditary permutability. No article title or link is supplied.
\end{itemize}
\end{document}
